\documentclass[aspectratio=169,11pt]{beamer}
\usepackage{../../../shared/theme/esmad_beamer_theme}
\usepackage{amsmath,amssymb}
\usepackage{booktabs}

\title[Measurement Error]{Measurement Error and Misclassification in Causal Inference}
\subtitle{When Variables Are Observed Imperfectly}
\date{\today}

\begin{document}

\begin{frame}
\titlepage
\end{frame}

\begin{frame}{Learning Goals}
\begin{itemize}
\item distinguish measurement error from missingness;
\item explain regression dilution under classical covariate error;
\item understand treatment and outcome misclassification;
\item distinguish differential from nondifferential error;
\item recognize residual confounding from noisy covariates;
\item use validation data, calibration, SIMEX, and sensitivity analysis.
\end{itemize}
\end{frame}

\begin{frame}{Outline}
\tableofcontents
\end{frame}

\section{Measurement Error as a Causal Problem}

\begin{frame}{True and Observed Variables}
Suppose the true variable is \(X\), but we observe
\[
X^\ast = X + U.
\]

The analysis conditions on or regresses on \(X^\ast\), not \(X\).
\end{frame}

\begin{frame}{Measurement Error Is Not Missing Data}
Missingness removes observations or values.

Measurement error retains a value, but the recorded value differs from the latent target.

\begin{alertblock}{Implication}
Bias can remain invisible because the dataset looks complete.
\end{alertblock}
\end{frame}

\section{Classical Measurement Error}

\begin{frame}{Classical Error Model}
A common model is
\[
X^\ast = X + U,
\]
with
\[
\E[U]=0,
\qquad
U \perp X.
\]

This is restrictive but useful for intuition.
\end{frame}

\begin{frame}{Regression Dilution}
In simple linear regression, classical error in an exposure tends to attenuate the slope toward zero.

If
\[
Y=\beta X+\varepsilon,
\]
but we regress on \(X^\ast\), then \(\widehat\beta\) is generally attenuated.
\end{frame}

\begin{frame}{Attenuation Is Not Universal}
The “bias toward zero” rule can fail when:
\begin{itemize}
\item multiple covariates are mismeasured;
\item the error is differential;
\item the outcome model is nonlinear;
\item the mismeasured variable is a confounder rather than the exposure.
\end{itemize}
\end{frame}

\section{Noisy Confounders}

\begin{frame}{Residual Confounding}
Suppose \(X\) is a confounder:
\[
X \rightarrow D,
\qquad
X \rightarrow Y.
\]

If only \(X^\ast\) is observed, conditioning on \(X^\ast\) may not block the confounding path fully.
\end{frame}

\begin{frame}{Measurement Quality Matters for Adjustment}
Even a correctly chosen adjustment set can fail if confounders are poorly measured.

\begin{alertblock}{Key lesson}
“Adjusted for age, income, severity...” is not enough; measurement reliability of those variables matters.
\end{alertblock}
\end{frame}

\section{Treatment Misclassification}

\begin{frame}{Observed vs True Treatment}
Let
\[
D^\ast
\]
be the recorded treatment and
\[
D
\]
the true treatment.

Misclassification occurs when
\[
P(D^\ast\neq D)>0.
\]
\end{frame}

\begin{frame}{Sensitivity and Specificity}
For binary treatment:
\[
Se=P(D^\ast=1\mid D=1),
\]
\[
Sp=P(D^\ast=0\mid D=0).
\]

These quantities characterize treatment classification quality.
\end{frame}

\begin{frame}{Nondifferential Treatment Misclassification}
If treatment misclassification is independent of outcome given true treatment and covariates, bias may often attenuate effects, but not always.

The direction depends on prevalence, sensitivity, specificity, and the estimand.
\end{frame}

\begin{frame}{Differential Treatment Misclassification}
If recording quality depends on outcome or outcome-related variables, bias can move in either direction.

Example:
\begin{itemize}
\item treatment adherence is self-reported more accurately by participants with better outcomes.
\end{itemize}
\end{frame}

\section{Outcome Misclassification}

\begin{frame}{Binary Outcome Misclassification}
Let \(Y^\ast\) be the measured outcome.

Sensitivity:
\[
P(Y^\ast=1\mid Y=1).
\]

Specificity:
\[
P(Y^\ast=0\mid Y=0).
\]
\end{frame}

\begin{frame}{Differential Outcome Measurement}
Outcome measurement may depend on treatment status.

Examples:
\begin{itemize}
\item treated patients receive more intensive monitoring;
\item treated firms are audited more often;
\item intervention participants report symptoms differently.
\end{itemize}

This can generate apparent treatment effects from detection alone.
\end{frame}

\section{Timing Error}

\begin{frame}{Treatment Timing Misclassification}
Suppose treatment actually begins at \(T_i\), but recorded time is
\[
T_i^\ast=T_i+\Delta_i.
\]

This can distort:
\begin{itemize}
\item event studies;
\item dynamic treatment effects;
\item interrupted time series;
\item anticipation analyses.
\end{itemize}
\end{frame}

\begin{frame}{Event-Time Smearing}
Timing error spreads a sharp causal response across adjacent event-time bins.

Observed event-study coefficients may appear delayed, attenuated, or falsely anticipatory.
\end{frame}

\section{Validation Data}

\begin{frame}{Validation Subsamples}
A validation sample contains both:
\begin{itemize}
\item the noisy measure;
\item a higher-quality reference measure.
\end{itemize}

It can estimate:
\begin{itemize}
\item reliability;
\item sensitivity and specificity;
\item measurement-error variance;
\item calibration functions.
\end{itemize}
\end{frame}

\begin{frame}{External vs Internal Validation}
Internal validation:
\begin{itemize}
\item occurs within the study sample.
\end{itemize}

External validation:
\begin{itemize}
\item uses a separate sample.
\end{itemize}

Transporting error parameters between populations requires its own assumptions.
\end{frame}

\section{Regression Calibration}

\begin{frame}{Regression Calibration}
Replace a noisy covariate with
\[
\E[X\mid X^\ast,Z],
\]
estimated from validation data or a measurement model.

This can reduce bias in some regression settings.
\end{frame}

\begin{frame}{Limits of Calibration}
Regression calibration may be inadequate for:
\begin{itemize}
\item strongly nonlinear outcome models;
\item binary misclassification;
\item differential measurement error;
\item complex causal structures.
\end{itemize}
\end{frame}

\section{SIMEX}

\begin{frame}{Simulation Extrapolation}
SIMEX:
\begin{enumerate}
\item adds extra simulated measurement error;
\item estimates the parameter at increasing error levels;
\item models the bias trajectory;
\item extrapolates back to the zero-error case.
\end{enumerate}
\end{frame}

\begin{frame}{SIMEX Assumptions}
SIMEX requires:
\begin{itemize}
\item a model for measurement-error variance;
\item a stable extrapolation function;
\item correct understanding of the error mechanism.
\end{itemize}
\end{frame}

\section{Quantitative Bias Analysis}

\begin{frame}{Bias Parameters}
For binary misclassification, specify plausible ranges for:
\begin{itemize}
\item sensitivity;
\item specificity;
\item differential error by treatment or outcome;
\item prevalence of the true exposure/outcome.
\end{itemize}
\end{frame}

\begin{frame}{Corrected Estimates}
Recompute the causal effect over plausible measurement-error scenarios.

\begin{block}{Goal}
Ask whether realistic error could explain, attenuate, or reverse the observed effect.
\end{block}
\end{frame}

\section{Measurement Error in Instruments and Mediators}

\begin{frame}{Instrument Measurement Error}
Noisy instruments can weaken relevance and increase finite-sample instability.

If measurement error changes the instrument-outcome relation directly, exclusion concerns may also arise.
\end{frame}

\begin{frame}{Mediator Measurement Error}
Mismeasured mediators can bias direct and indirect effects.

Because mediation is already assumption-heavy, poor mediator measurement compounds identification uncertainty.
\end{frame}

\section{Applied Reasoning}

\begin{frame}{Example: Regional Pricing}
Suppose the treatment is a regional pricing intervention, but implementation dates are recorded with error.

Possible consequences:
\begin{itemize}
\item event-study effects shift across time;
\item pre-trend diagnostics become contaminated;
\item treatment intensity is misclassified;
\item spillovers may be mistaken for measurement error.
\end{itemize}
\end{frame}

\begin{frame}{A Credible Measurement-Error Workflow}
\begin{enumerate}
\item Identify which variables may be measured with error.
\item Define the latent target and observed proxy.
\item Determine whether error is classical, differential, or misclassification.
\item Use validation data when available.
\item Quantify reliability, sensitivity, and specificity.
\item Apply calibration or measurement models where justified.
\item Run quantitative bias or SIMEX-style sensitivity analyses.
\item Report corrected and uncorrected estimates.
\end{enumerate}
\end{frame}

\begin{frame}{Common Failure Modes}
\begin{itemize}
\item assuming all measurement error attenuates toward zero;
\item ignoring noisy confounders;
\item treating self-reported treatment as exact;
\item overlooking differential outcome detection;
\item applying calibration without a valid measurement model;
\item using validation parameters from a non-comparable population.
\end{itemize}
\end{frame}

\begin{frame}{Synthesis: Measurement Is Part of Identification}
\begin{itemize}
\item Measurement error can affect treatment, outcomes, confounders, mediators, and timing.
\item Classical attenuation is only one special case.
\item Noisy confounders leave residual confounding.
\item Differential misclassification can bias effects in either direction.
\item Validation data can identify key error parameters.
\item Calibration, SIMEX, and bias analysis require explicit error assumptions.
\item Measurement quality belongs in the causal argument, not only the data-cleaning appendix.
\end{itemize}
\end{frame}

\begin{frame}{Selected References}
\small
\begin{itemize}
\item Carroll, R. J. et al. (2006). \textit{Measurement Error in Nonlinear Models}, 2nd ed. Chapman \& Hall/CRC.
\item Stefanski, L. A., and Cook, J. R. (1995). Simulation-Extrapolation: The Measurement Error Jackknife. \textit{JASA}, 90(432), 1247--1256.
\item Gustafson, P. (2004). \textit{Measurement Error and Misclassification in Statistics and Epidemiology}. Chapman \& Hall/CRC.
\item Greenland, S., and Lash, T. L. (2008). Bias Analysis. In \textit{Modern Epidemiology}, 3rd ed.
\item Keogh, R. H., and White, I. R. (2014). A Toolkit for Measurement Error Correction. \textit{Stata Journal}, 14(2), 301--325.
\end{itemize}
\end{frame}

\contactslide

\end{document}
